\section{欧氏关联函数的长程渐近：推广}\label{sec:yzliu-decay}

第~22~章的 Fourier 步骤依赖类空关联函数的指数衰减（\cref{thm:lamet-exp-decay}）。Ji、Liu 与 Su\cite{Ji:2026vir}（预印本）把这一结果推广到有单粒子中间态的道，并在标量理论中给出微扰论各阶的证明。

\begin{proposition}[标量场论中的谱表示与指数衰减（微扰论各阶） \tagP]\label{prop:yzliu-decay-extensions}
在 $\varphi^3$ 与 $\varphi^4$ 标量理论中，微扰论各阶，任意外动量下的前向关联函数有谱表示
\begin{equation}\label{eq:yzliu-decay-spectral}
f(\vec z,\vec p)=2Z\cos(\vec z\cdot\vec p)+\int_{-1}^1\dd x\int_{s(x)}^\infty\dd s\,\rho(s,x)\,\ee^{\ii x\vec z\cdot\vec p}K_0(s|\vec z|),
\end{equation}
谱支撑 $x\in(-1,1)$（“a result of locality”），谱隙 $s_0=\min_xs(x)>0$；因此它在大类空间隔处指数衰减（由参数表示证明\cite[第~5.3~节]{Ji:2026vir}）。两圈 $\varphi^3$ 计算检验了渐近形式。（本命题只涉及标量理论。）
% uses: none
\end{proposition}
% Writer note (W26): eq. (5.33) of Ji:2026vir checked on the PDF (v1). In its eqs. (5.21)-(5.23) the third exponent lacks |z| and
% (5.22) carries m^2 instead of m in front of the square root (dimensional slips); only the verbal content is printed.

\paragraph{地位与缺口。}[P]（标量理论，微扰论；预印本\cite[第~1--2、5、6~节]{Ji:2026vir}）。同一文献经形状因子的谱表示及其“局域性完备化”，把\cref{thm:lamet-exp-decay}推广到有单粒子中间态的道：衰减率为 $M$、$\sqrt{s_{\mathrm{th}}}\,m$ 或由 $s_{\mathrm{th}}$、$m$、$M$ 决定的 $m_c$（$m$、$M$ 为外态与最轻中间单粒子态的质量，$s_{\mathrm{th}}m^2$ 为单粒子形状因子在类时道的阈值，与上式的谱隙 $s_0$ 不是同一个量）\cite[式~(5.19)--(5.23)]{Ji:2026vir}，局域性消去了末态形状因子的奇点。这一推广的地位与\cref{thm:lamet-exp-decay}相同——在公理下的蕴含，其内容等级经谱假设为 [C]——不另行定级。关联：第~22~章的 Fourier 步骤（\cref{lem:lamet-strip}、\cref{prop:lamet-truncation}）与正问题／逆问题之争。
